\section{Completion of the realization construction}
\label{sec:realization}

We now prove Proposition~\ref{prop:realization}.  The immersion of
Proposition~\ref{prop:immersion} has been placed in a proper uncurved
cyclic category by Proposition~\ref{prop:curvature}, and
Proposition~\ref{prop:directed-comparison}
identifies the directed operations on every finite collection of
embedded copies with those on the corresponding mirror tests.
We first assemble these finite comparisons.  Cofinal copies will then
recover the full mirror category, to which the immersed object admits
an orthogonal projection.

\subsection{Copies and a common coefficient field}

Retain the ordered Hamiltonian copies chosen in
Section~\ref{sec:curvature}.  Each of the finitely many tests
$K_{\ell,M}$ occurs cofinally: after every label there is a later copy
of each test.
Finite sets in the exhaustion will always include the original tests
and at least three copies of each test.  Denote
the immersed object by $L$, with its grading and spin data understood.

\begin{lemma}[Invertibility of copy arrows]
\label{realization:copies}
In each finite immersed category containing at least three copies of
each test, every nonzero degree-zero cohomology class between distinct
copies of the same test is invertible.
\end{lemma}

\begin{proof}
The strict pair comparison gives a one-dimensional degree-zero
cohomology space between copies.  Given three distinct copies, choose
their directed order in any desired way.
Proposition~\ref{prop:directed-comparison}, ordinary HMS, and
Hamiltonian-copy invariance in each factor show that the product of two
nonzero composable degree-zero classes between those copies is nonzero.

Let $u:A\to B$ and $v:B\to A$ be such classes.  Choose a third copy $C$
and a nonzero class $w:A\to C$.  If $vu=0$, associativity would give
$0=w(vu)=(wv)u$, contrary to the nonzero-product assertion on distinct
copies.  Thus $vu\ne0$, and similarly $uv\ne0$.  The self degree-zero
cohomology has dimension at most one, because its canonical chain space
in that degree is $H^0(K_{\ell,M})$.  The cohomological units therefore
make $vu$ and $uv$ nonzero multiples of the identities.  Rescaling one
arrow gives inverse classes.
\end{proof}

The ordinary mirrors of all copies can be fixed once and for all.
Indeed, first choose the finite ordinary factor categories on the
original tests and their specified images $Q_\ell,N_M$.  For each larger
finite system, compare ordinary choices after including these original
labels.  Hamiltonian copies are isomorphic to them in the ordinary
categories.  The finite comparison can consequently use the same
mirror objects, with repeated labels.  No analytic functor on an
infinite collection is needed.

The finite cyclic categories and their comparisons need not be
compatible as the label set increases.  We assemble them by a field
ultraproduct: the canonical complexes bound each Hom dimension, and
every equation involving fixed labels is a finite matrix identity.
These two facts allow all the required structures to pass to one
coefficient field in the next lemma.

\begin{lemma}[Simultaneous finite comparisons]
\label{realization:field}
Over an extension field $\Bbbk$ of characteristic zero and cardinality
at most the continuum, there is a proper uncurved cyclic
$A_\infty$ category containing $L$ and all the test copies.  Its
directed subcategory on the test copies admits an $A_\infty$ functor to fixed mirror
models, inducing isomorphisms on the cohomology of every strict arrow.
The degree-two bound for $L$ is
\begin{equation}
 \dim_{\Bbbk}\Hom^2(L,L)
 \le b_2(D)-\dim\mathcal Z.
 \label{realization:ambient-bound}
\end{equation}
The category has cohomological units and a nondegenerate cyclic pairing
of dimension $n$.
\end{lemma}

\begin{proof}
Exhaust the labels by finite systems.  For each one, take the uncurved
$A_\infty$ category of Proposition~\ref{prop:curvature}, with its bound
\eqref{curv:degree-two-bound}, and the
directed comparison of Proposition~\ref{prop:directed-comparison},
extending its coefficient field when necessary.  Pass to minimal
models.  The immersed models can be chosen cyclic: a finite complex
with a nondegenerate compatible pairing admits a splitting and a
contraction compatible with that pairing.  Transfer the directed
functors to these models and to the fixed mirror models on the repeated
labels.

For any fixed pair of labels, the dimensions and degrees of the minimal
Hom spaces are bounded by the ranks of the corresponding finite
canonical complexes.  Choose a nonprincipal ultrafilter on the
exhaustion.  For each fixed Hom, its graded ranks are constant on a set
in the ultrafilter.  Use bases of those ranks and take the ordinary
field ultraproduct of the coefficient fields.  Every fixed
$A_\infty$ equation and functor equation is a finite matrix identity,
so passes to this field.  The same applies to units, cyclicity,
nondegeneracy of the pairings, invertibility of the strict-arrow
cohomology maps and of the copy arrows, and
\eqref{realization:ambient-bound}.
The fixed mirror models undergo ordinary scalar extension along the
diagonal embedding of their common ground field.

The finite canonical ranks show that the resulting category is proper.
All fields in the finite constructions can be taken to have cardinality
at most the continuum, and a countable ultraproduct has the same bound.
There is therefore an abstract embedding of the resulting field into
$\C$.  The construction has used ordinary finite-object HMS before
taking this ultraproduct; it makes no assertion of analytic mirror
symmetry over an ultraproduct field.
\end{proof}

\subsection{Recovering full morphisms from directed copies}

We explain why the directed comparison determines the enhanced full
subcategory of copies.  The two receiving categories are the full
subcategory of test copies in the immersed category and the full
subcategory of their specified mirror images, with repeated labels.
Replace both receiving categories by
strictly unital models, transport the functors on strict words, and
adjoin formal units to the directed category.  Write $\mathcal D$ for
this unital directed category, and let $\mathcal C$ denote either full
category of copies receiving it.  The receiving functor is a
quasi-isomorphism on strict arrows, and copy arrows become isomorphisms
in cohomology by Lemma~\ref{realization:copies} and its ordinary mirror
counterpart.

For a label $j$, write $h_j^{\mathcal D}=\Hom_{\mathcal D}(-,j)$ for its
right Yoneda module.  Choose successive later copies
$j_0,j_1,\ldots$ of $j$ and degree-zero copy arrows between them.
Restriction of the full representable admits a natural comparison
\begin{equation}
 \mathop{\operatorname{hocolim}}_k h_{j_k}^{\mathcal D}
 \longrightarrow
 \Hom_{\mathcal C}(-,j)|_{\mathcal D}.
 \label{realization:telescope}
\end{equation}
To construct it, first take the same telescope of restricted full
representables.  All its arrows are equivalences, so its homotopy
colimit is the right-hand side.  The Yoneda comparisons from directed
representables to restricted full ones commute with the arrows up to
the natural homotopies and induce \eqref{realization:telescope}.
At a fixed argument label $i$, every sufficiently late $j_k$ satisfies
$i<j_k$, where the comparison on Hom is a quasi-isomorphism.
Thus \eqref{realization:telescope} is a quasi-isomorphism on every
argument.

Taking Hom from this telescope to another restricted full
representable gives the homotopy limit of its values at $j_k$.
The transition maps in that homotopy limit are equivalences, because
the full-category copy arrows are invertible.  It follows that
restriction of Yoneda is fully faithful on the full copy category.
This argument applies to higher $A_\infty$ functors as well, using
their natural Yoneda comparisons.

On both sides the restricted representable of each original test is
therefore the same telescope in the derived module category of
$\mathcal D$.  The enhanced split-closed triangulated hulls of the two
full categories consequently agree.  By
Lemma~\ref{setup:graph-generation} and sphere generation, their mirror
hull is
\begin{equation}
 \mathcal A\simeq\Perf(B\times E^g).
 \label{realization:mirror-hull}
\end{equation}
Here we use ordinary external-product generation for smooth projective
varieties: external products of generators of the two factors generate
the perfect category of their product.

\subsection{Projection and the Euler pairings}

Let $\mathcal F$ be the split-closed triangulated hull of the immersed
category including $L$.  Its nondegenerate cyclic pairing makes it
proper with Serre functor $[n]$; the pairing extends to twisted complexes
and then to retracts.  The full subcategory $\mathcal A$ in
\eqref{realization:mirror-hull} is saturated
\cite{BondalVanDenBergh}.  The module
$\Hom_{\mathcal F}(-,L)|_{\mathcal A}$ has finite-dimensional total
cohomology at every object, by properness.  Saturation therefore
represents it by an object
$\mathcal E\in\mathcal A$.  The representing map fits into a triangle
\[
 \mathcal E\longrightarrow L\longrightarrow R\longrightarrow
 \mathcal E[1],
 \qquad \Hom^*(\mathcal A,R)=0.
\]
Serre duality in $\mathcal F$ gives
$\Hom^*(R,\mathcal A)=0$ as well.  The connecting map vanishes, and the
triangle splits.  In particular $\mathcal E$ is a summand of $L$, so
\[
 \dim\Ext^2(\mathcal E,\mathcal E)
 \le\dim\Hom^2_{\mathcal F}(L,L)
 \le b_2(D)-\dim\mathcal Z.
\]
After the abstract field embedding, this proves
\eqref{setup:ext-bound} on the complex mirror.

For every test object, the orthogonal summand $R$ has zero Euler
pairing.  The Euler pairing with $L$ is computed by its transverse
Floer complex, hence equals the oriented intersection pairing with
the immersed fundamental class, up to the common sign fixed by the
grading convention.  Equation~\eqref{topo:class} says
$\PD\,i_*[L]=\nu a$ with $\nu>0$.  It follows that all the test Euler
pairings equal the pairings with $a$ times one common nonzero rational
scalar.  This is \eqref{setup:euler} and completes the proof of
Proposition~\ref{prop:realization}.
